\chapter{统一多任务模型设计与训练优化}

\section{模型总体结构}

本文统一多任务模型定义在 \code{src/model/unified_model.py} 中，命名为 \code{UnifiedMultiTaskModel}。模型整体结构采用“YOLO 共享主干 + 两个附加头”的形式，即：

\begin{enumerate}[leftmargin=2.5em]
\item 使用 YOLO 主干和颈部完成多尺度特征提取；
\item 保留 YOLO 检测头用于车辆检测；
\item 在高层共享特征图上增加雾分类头；
\item 在同一共享特征图上增加 beta 回归头。
\end{enumerate}

这种结构的核心思想是：车辆检测需要局部多尺度感知，而天气识别和雾浓度估计更偏全局语义；因此，将这三项任务放在同一主干上，可以在兼顾速度的同时利用高层共享特征实现协同学习。相比使用三套独立模型，这种设计具有参数更少、推理链路更短、部署更容易的优势。

\begin{figure}[htbp]
  \centering
  \includegraphics[width=0.92\textwidth]{fig_03_model_structure.png}
  \caption{统一多任务模型结构图}
  \label{fig:model_structure}
\end{figure}

\section{单类车辆检测头适配}

本文场景主要关注道路车辆，而不需要保留通用目标检测中的 80 类全部语义。因此，系统在配置中明确设定 \texttt{NUM\_DET\_CLASSES = 1}，类别名称为 \texttt{["vehicle"]}。然而，基础权重 \texttt{yolo11n.pt} 原本是在多类别数据上训练得到的，其检测头输出维度与单类任务不一致。为解决这一问题，本文在模型初始化阶段检测基础模型最后一层的类别数；若与当前任务不一致，则重建单类 \texttt{DetectionModel}，并从原始权重中选择形状兼容的参数进行非严格加载。

这种适配策略十分关键。若直接以严格模式加载多类权重到单类模型，程序会因 shape 不匹配失败；若完全放弃预训练权重，又会失去 backbone 和 neck 中已学到的大量通用视觉特征。本文采取的折中方案是：保留能复用的主干参数，仅重置类别不兼容的检测头部分。这样既能保留迁移学习优势，又能让模型适配当前道路监控单类检测任务。

\section{雾分类头设计}

雾分类头的输入来自 YOLO 高层共享特征，而非直接使用原图像。其结构为：

\begin{center}
\small
\code{AdaptiveAvgPool2d -> Flatten -> Linear(256)}\\
\code{-> ReLU -> Dropout -> Linear(num_fog_classes)}
\end{center}

之所以采用全局平均池化，是因为雾类型识别更关注整幅画面的整体退化模式，而不是某个局部区域的细节。通过池化操作，模型可将空间特征压缩为全局描述向量，再通过全连接层完成三分类。当前雾类别定义为：

\begin{itemize}[leftmargin=2em]
\item \texttt{clear}
\item \texttt{uniform}
\item \texttt{patchy}
\end{itemize}

这三类设置兼顾了业务可理解性与训练可实现性。\texttt{clear} 表示无雾，\texttt{uniform} 表示均匀低能见度，\texttt{patchy} 则用于表示局地非均匀团雾。在实际管理中，\texttt{patchy} 比 \texttt{uniform} 更值得关注，因为它更容易诱发突发性风险。

\section{beta 回归头设计}

回归头同样基于共享高层特征设计，其结构为：

\texttt{AdaptiveAvgPool2d -> Flatten -> Linear(128) -> ReLU -> Linear(1) -> Sigmoid}

其中最终使用 \texttt{Sigmoid} 将输出压缩到 \texttt{[0,1]} 范围，再在训练与推理阶段按 \texttt{Config.BETA\_MAX} 做缩放。这意味着模型回归的并不是原始物理量纲下的绝对散射系数，而是相对于预设上界的归一化值。当前实现中，回归目标在 clear 样本上被显式置为 0，在有雾样本上由在线造雾过程给出真实采样 beta。因此，该回归头可以被视为学习“当前画面相对于训练雾强度上限的退化程度”。

beta 回归的意义并不只是提供一个数字。对于系统应用而言，连续 beta 值可以进一步用于显示平滑、动态阈值控制、预警等级划分和小地图可视化，因此它是连接模型输出与业务联动的重要桥梁。

\section{多任务前向传播与共享特征提取}

与许多简单封装直接调用 \code{self.yolo(x)} 的方式不同，本文在 \code{forward()} 中手动遍历 \code{self.yolo.model} 各层前向传播，并尝试捕获 SPPF 或检测头前最后一张高层 4D 特征图作为共享特征。这一实现看似复杂，但非常必要。因为若只调用高层封装接口，往往难以在推理结构变化和训练/验证模式差异下稳定获得可供附加头使用的高层特征。

本文采用的共享特征提取策略是：优先捕获 \texttt{SPPF} 层输出；若当前 YOLO 架构中无法明确识别 \texttt{SPPF}，则退化为检测头前最后一张 4D 特征图。这种设计提升了模型对不同 YOLO 规模和结构变体的适应性，也避免了把通道数写死在代码里。

\section{损失函数与联合优化}

训练阶段，检测、分类和回归三项任务同时纳入优化目标。具体而言：

\begin{itemize}[leftmargin=2em]
\item 检测损失直接调用 Ultralytics 提供的 detection loss；
\item 雾分类损失使用 \texttt{CrossEntropyLoss}；
\item beta 回归损失使用 \texttt{MSELoss}。
\end{itemize}

总损失写为：

\[
L = L_{det} + L_{cls} + L_{reg}
\]

虽然当前实现中三项任务权重均为 1.0，但这并不意味着三项损失在数值尺度上天然相等。实际训练中，检测损失往往明显高于雾分类与回归损失，这也是当前日志中“det 占主导，fog\_cls 与 fog\_reg 数值更小”的原因。对此，本文采取两种应对方式：一是通过共享高层特征让天气任务充分利用检测主干已有的表达能力，减轻天气头从零学习的压力；二是在实验设计中预留损失权重调整空间，以便后续根据收敛曲线进一步平衡不同任务。

\section{训练稳定性与工程保护机制}

相比只关注最终精度的论文实现，本文训练脚本在工程稳定性方面做了较多保护设计。

首先，训练过程支持自动续训。系统会在 checkpoint 目录中查找最新权重，若存在则恢复模型、优化器、调度器和 AMP scaler 状态，并从下一个 epoch 继续训练。这对于长时间训练、多次中断和实验迭代都十分重要。

其次，训练过程支持缺失深度缓存自动补算。训练开始前会统计训练集和验证集各自缺少多少深度文件，若有缺失则自动调用预计算流程补齐，避免训练到验证阶段才因文件缺失报错。

再次，脚本支持 AMP 混合精度训练，并监控非有限梯度比例。当前代码中实现了对 \texttt{NaN/Inf} 梯度的检测、告警、阈值判定以及必要时自动关闭 AMP 的恢复策略。这种设计在大 batch、长时间训练或量化感知训练切换中尤为重要。

最后，系统保留了 QAT 与 INT8 转换路径。虽然当前毕业设计阶段更关注功能打通与方法验证，但在 \texttt{src/train.py} 中已经实现了 QAT 准备、observer 修复和 INT8 模型转换流程，为后续边缘部署提供了明确方向。

\section{当前训练状态与续训策略}

在本文撰写时，项目工作区中已经存在一版训练权重和 checkpoint。现有 checkpoint 对应的训练记录显示，系统此前至少完成了首个 epoch 的训练，并保存了 \code{unified_model.pt}、\code{unified_model_best.pt} 以及 \code{checkpoint_epoch_0001.pt}。从工程验证看，这意味着统一模型主链已经能够完成参数更新、权重保存和推理调用；从效果层面看，早期权重在真实下载雾视频上的雾估计已可输出数值，但车辆检测稳定性仍不如直接使用预训练 YOLO11n 的结果，因此本文在系统演示阶段补充实现了混合推理策略。

为了进一步提升雾分类和 beta 回归可靠性，当前系统已在后台执行一轮“短程快速微调”：继续沿用现有 checkpoint，但对每个 epoch 的最大训练 batch 数和验证 batch 数做限制，从而在可接受时间内获得更新权重。这一策略的意义不在于替代完整训练，而在于为论文初稿撰写阶段提供更接近最终状态的参数版本，并为后续实验章节补充完整指标争取时间。
